%% TRANSCRIBED (PDF 20-24)

\sechead{1}{The notion of an \textit{l}-adic representation}%
  {THE NOTION OF AN \textit{l}-ADIC REPRESENTATION}

\subsechead{1.1}{Definition}

Let $K$ be a field, and let $K_s$ be a separable algebraic closure of
$K$.  Let $G = \Gal(K_s/K)$ be the Galois group of the extension
$K_s/K$.  The group $G$, with the Krull topology, is compact and totally
disconnected.  Let $\ell$ be a prime number, and let $V$ be a
finite-dimensional vector space over the field $\Q_\ell$ of
\textit{l}-adic numbers.  The full linear group $\Aut(V)$ is an
\textit{l}-adic Lie group, its topology being induced by the natural
topology of $\End(V)$; if $n = \dim(V)$, we have
$\Aut(V) \simeq \GL(n, \Q_\ell)$.

\begin{dfn}
An \textit{l}-adic representation of $G$ (or, by abuse of language, of
$K$) is a continuous homomorphism $\rho : G \longrightarrow \Aut(V)$.
\end{dfn}

\begin{remarks}
\begin{nlist}
\item A \emph{lattice} of $V$ is a sub-$\Z_\ell$-module $T$ which is
free of finite rank, and generate $V$ over $\Q_\ell$, so that $V$ can be
identified with $T \tens_{\Z_\ell} \Q_\ell$.  Notice that
\emph{there exists a lattice of $V$ which is stable under} $G$.
This follows from the fact that $G$ is compact.

Indeed, let $L$ be any lattice of $V$, and let $H$ be the set of
elements $g \in G$ such that $\rho(g)L = L$.  This is an open subgroup
of $G$, and $G/H$ is finite.  The lattice $T$ generated by the lattices
$\rho(g)L$, $g \in G/H$, is stable under $G$.

Notice that $L$ may be identified with the projective limit of the free
$(\Z/\ell^m\Z)$-modules $T/\ell^m T$, on which $G$ acts; the vector
space $V$ may be reconstructed from $T$ by
$V = T \tens_{\Z_\ell} \Q_\ell$.

\item If $\rho$ is an \textit{l}-adic representation of $G$, the group
$G_\rho = \Ima(\rho)$ is a closed subgroup of $\Aut(V)$, and hence, by
the \textit{l}-adic analogue of Cartan's theorem
(cf.~\cite{serre-liealg}, LG, p.~5-42) $G_\rho$ is itself an
\emph{\textit{l}-adic Lie group}.  Its Lie algebra
$\mathfrak{g}_\rho = \mathrm{Lie}(G_\rho)$ is a subalgebra of
$\End(V) = \mathrm{Lie}(\Aut(V))$.  The Lie algebra $\mathfrak{g}_\rho$
is easily seen to be invariant under extensions of finite type of the
ground field $K$ (cf.~\cite{serre-congr}, 1.2).
\end{nlist}
\end{remarks}

\begin{exercises}
\begin{nlist}
\item Let $V$ be a vector space of dimension $2$ over a field $k$ and
let $H$ be a subgroup of $\Aut(V)$.  Assume that $\det(1-h) = 0$ for all
$h \in H$.  Show the existence of a basis of $V$ with respect to which
$H$ is contained either in the subgroup
$\bigl(\begin{smallmatrix} 1 & * \\ 0 & * \end{smallmatrix}\bigr)$ or in
the subgroup
$\bigl(\begin{smallmatrix} 1 & 0 \\ * & * \end{smallmatrix}\bigr)$ of
$\Aut(V)$.

\item Let $\rho : G \longrightarrow \Aut(V_\ell)$ be an \textit{l}-adic
representation of $G$, where $V_\ell$ is a $\Q_\ell$-vector space of
dimension $2$.  Assume $\det(1-\rho(s)) \equiv 0 \bmod.\ \ell$ for all
$s \in G$.  Let $T$ be a lattice of $V_\ell$ stable by $G$.  Show the
existence of a lattice $T'$ of $V_\ell$ with the following two
properties.
\begin{enumerate}[label=\alph*),leftmargin=2.2em,itemsep=0.15ex,topsep=0.4ex]
\item $T'$ is stable by $G$
\item Either $T'$ is a sublattice of index $\ell$ of $T$ and $G$ acts
trivially on $T/T'$ or $T$ is a sublattice of index $\ell$ of $T'$ and
$G$ acts trivially on $T'/T$.
\end{enumerate}
(Apply exercise~1) above to $k = \F_\ell$ and $V = T/\ell T$.)

\item Let $\rho$ be a semi-simple \textit{l}-adic representation of $G$
and let $U$ be an invariant subgroup of $G$.  Assume that, for all
$x \in U$, $\rho(x)$ is unipotent (all its eigenvalues are equal to
$1$).  Show that $\rho(x) = 1$ for all $x \in U$.  (Show that the
restriction of $\rho$ to $U$ is semi-simple and use Kolchin's theorem to
bring it to triangular form.)

\item Let $\rho : G \longrightarrow \Aut(V_\ell)$ be an \textit{l}-adic
representation of $G$, and $T$ a lattice of $V_\ell$ stable under $G$.
Show the equivalence of the following properties:
\begin{enumerate}[label=\alph*),leftmargin=2.2em,itemsep=0.15ex,topsep=0.4ex]
\item The representation of $G$ in the $\F_\ell$-vector space $T/\ell T$
is irreducible.
\item The only lattices of $V_\ell$ stable under $G$ are the $\ell^n T$,
with $n \in \Z$.
\end{enumerate}
\end{nlist}
\end{exercises}

\subsechead{1.2}{Examples}

\begin{enumerate}[label=\arabic*.,leftmargin=2.2em,itemsep=0.15ex,topsep=0.4ex]
\item \emph{Roots of unity.}  Let $\ell \neq \mathrm{char}(K)$.  The
group $G = \Gal(K_s/K)$ acts on the group $\mu_m$ of $\ell^m$-th roots
of unity, and hence also on $T_\ell(\mu) = \varprojlim\ \mu_m$.  The
$\Q_\ell$-vector space $V_\ell(\mu) = T_\ell(\mu) \tens_{\Z_\ell}
\Q_\ell$ is of dimension $1$, and the homomorphism
$\chi_\ell : G \longrightarrow \Aut(V_\ell) = \Q_\ell^*$ defined by the
action of $G$ on $V_\ell$ is a $1$-dimensional \textit{l}-adic
representation of $G$.  The character $\chi_\ell$ takes its values in
the group of units $U_\ell$ of $\Z_\ell$; by definition
\[
g(z) = z^{\chi_\ell(g)} \quad\text{if } g \in G, \quad z^{\ell^m} = 1.
\]

\item \emph{Elliptic curves.}  Let $\ell \neq \mathrm{char}(K)$.  Let
$E$ be an elliptic curve defined over $K$ with a given rational point
$0$.  One knows that there is a unique structure of group variety on $E$
with $0$ as neutral element.  Let $E_m$ be the kernel of multiplication
by $\ell^m$ in $E(K_s)$, and let
\[
T_\ell(E) = \varprojlim\ E_m, \quad V_\ell(E) = T_\ell(E) \tens_{\Z_\ell} \Q_\ell.
\]
The Tate module $T_\ell(E)$ is a free $\Z_\ell$-module on which
$G = \Gal(K_s/K)$ acts (cf.~\cite{lang-abelian}, chap.~VII).  The
corresponding homomorphism $\pi_\ell : G \longrightarrow \Aut(V_\ell(E))$
is an \textit{l}-adic representation of $G$.  The group
$G_\ell = \Ima(\pi_\ell)$ is a closed subgroup of $\Aut(T_\ell(E))$, a
$4$-dimensional Lie group isomorphic to $\GL(2, \Z_\ell)$.  (In
chapter~\ref{chap:IV}, we will determine the Lie algebra of $G_\ell$,
under the assumption that $K$. is a number field.)
%% NOTE: the stray period after "$K$" above is printed in the original
%% (I-4, both text layer and image); kept as printed.

Since we can identify $E$ with its dual (in the sense of the duality of
abelian varieties) the symbol $(x, y)$ (cf.~\cite{lang-abelian},
\emph{loc.~cit.}) defines canonical isomorphisms
\[
\Lambda^2 T_\ell(E) = T_\ell(\mu), \quad \Lambda^2 V_\ell(E) = V_\ell(\mu).
\]
Hence $\det(\pi_\ell)$ is the character $\chi_\ell$ defined in
example~1.

\item \emph{Abelian varieties.}  Let $A$ be an abelian variety over $K$
of dimension $d$.  If $\ell \neq \mathrm{char}(K)$, we define
$T_\ell(A)$, $V_\ell(A)$ in the same way as in example~2.  The group
$T_\ell(A)$ is a free $\Z_\ell$-module of rank $2d$
(cf.~\cite{lang-abelian}, \emph{loc.~cit.}) on which
$G = \Gal(K_s/K)$ acts.

\item \emph{Cohomology representations.}  Let $X$ be an algebraic
variety defined over the field $K$, and let $X_s = X \times_K K_s$ be
the corresponding variety over $K_s$.  Let $\ell \neq
\mathrm{char}(K)$, and let $i$ be an integer.  Using the étale
cohomology of Artin-Grothendieck~\cite{artin-grothendieck-sga} we let
\[
H^i(X_s, \Z_\ell) = \varprojlim\ H^i((X_s)_{\text{ét}}, \Z/\ell^n\Z),
\]
\[
H^i_\ell(X_s) = H^i(X_s, \Z_\ell) \tens_{\Z_\ell} \Q_\ell.
\]
The group $H^i_\ell(X_s)$ is a vector space over $\Q_\ell$ on which
$G = \Gal(K_s/K)$ acts (via the action of $G$ on $X_s$).  It is finite
dimensional, at least if $\mathrm{char}(K) = 0$ or if $X$ is proper.  We
thus get an \textit{l}-adic representation of $G$ associated to
$H^i_\ell(X_s)$; by taking duals we also get homology \textit{l}-adic
representations.  Examples~1, 2, 3 are particular cases of homology
\textit{l}-adic representations where $i = 1$ and $X$ is respectively
the multiplicative group $\Gm$, the elliptic curve $E$, and the abelian
variety $A$.
\end{enumerate}

\runhead{Exercise}
\begin{alist}
\item Show that there is an elliptic curve $E$, defined over
$K_o = \Q(T)$, with $j$-invariant equal to $T$.
%% NOTE: the original prints the subscript of $K$ as a small round
%% lowercase "o" (typewriter nought), clearly distinct from the narrow
%% oval digit "0" of the same font (cf. "char(K) = 0", same page);
%% transcribed literally as $K_o$.
\item Show that for such a curve, over $K = \C(T)$, one has
$G_\ell = \SL(T_\ell(E))$ (cf.~Igusa~\cite{igusa} for an algebraic
proof).
\item Using (b), show that, over $K_o$, we have
$G_\ell = \GL(T_\ell(E))$.
\item Show that for any closed subgroup $H$ of $\GL(2, \Z_\ell)$ there
is an elliptic curve (defined over some field) for which $G_\ell = H$.
\end{alist}
